
\begin{figure}[t]
    \centering
\begin{lstlisting}[style=mypython]

class HasReasoning(dspy.Signature):
    """Determine if the model response contains reasoning steps."""
    model_response: str = dspy.InputField()
    has_reasoning: bool = dspy.OutputField()

class HasReasoningOutline(dspy.Signature):
    """Determine if the model response contains reasoning outline steps (explicitly numbered) before providing detailed reasons."""
    model_response: str = dspy.InputField()
    has_reasoning_outline: bool = dspy.OutputField()

class HasSkipThoughts(dspy.Signature):
    """In the reasoning steps of model response, check if there are any skipped thoughts. 
    Example of thought skipping: 
        model response: ``Good question. Let's think about steps.\n1. Identify candidates; 2. Compare their weights.\nHydrogen and Carbon are possible candidates. The answer is B.'' 
        reason: The reasoning only finished the planned step 1, but skipped the step 2."""
    model_response: str = dspy.InputField()
    has_skip_thoughts: bool = dspy.OutputField()

class FactualErrorInReasoning(dspy.Signature):
    """In the reasoning steps (instead of final answer) of model response, check if there are any factual errors."""
    query_to_model: str = dspy.InputField()
    model_response: str = dspy.InputField()
    identified_factual_errors: List[str] = dspy.OutputField()
    has_factual_error: bool = dspy.OutputField()

class HasWrongLogic(dspy.Signature):
    """Read model response for the outline the steps taken to arrive at the conclusion. Check if there are any wrong logic errors in the outline."""
    query_to_model: str = dspy.InputField()
    model_response: str = dspy.InputField()
    identified_logic_errors: List[str] = dspy.OutputField()
    has_logic_error: bool = dspy.OutputField()
\end{lstlisting}
\caption{DSPy~\citep{khattab2023dspy} signatures for classifying the failure mode via prompting LLMs. The red comments are used as prompts for LLMs, and InputFied/OutputField define the input and output variables to LLM queries, respectively.}
\label{fig:failure_clf_prompt}
\end{figure}

% \begin{figure}[t]
%     \centering
%     % colback=gray!5!white,colframe=black!50!gray
%     \begin{boxK}[colback=royalblue!5!white, colframe=royalblue!75!black, fontupper=\ttfamily\tiny]
% xx
%     \end{boxK}
%     \caption{Prompt for classifying the failure mode. The failure mode is a multi-choice categorization. \junyuan{to update}}
%     \label{fig:failure_clf_prompt}
% \end{figure}

% \begin{figure}[t]
%     \centering
%     % colback=gray!5!white,colframe=black!50!gray
%     \begin{boxK}[colback=royalblue!5!white, colframe=royalblue!75!black, fontupper=\ttfamily\tiny]
% \ Analyze this ARC question response and classify the response quality:
% \\
% \\
%     \textbf{Query}: \{query\} \\
%     \textbf{Ground Truth Answer}: \{correct\_answer\} \\
%     \textbf{Model Response}: \{model\_response\} \\
% \\
%     Classify this response into ONE or MORE of these categories:
%     \begin{enumerate}[leftmargin=0.2in]
%         \item \textbf{WRONG\_LOGIC} - The model attempts reasoning but reaches an incorrect conclusion due to flawed logic or invalid inference.
%         \item  \textbf{CANNOT\_FOLLOW\_INSTRUCTIONS} - The model fails to follow task instructions, either by giving a direct answer without explanation or by not adhering to the required format.
%         \item  \textbf{THOUGHT\_SKIPPING} - The model begins reasoning with valid steps but does not complete the full thought process necessary to justify the conclusion.
%         \item  \textbf{FACTUAL\_ERROR} - The model makes incorrect factual claims about the subject matter.
%         \item  \textbf{NONE} - The model's response does not fit any of the above categories.
%     \end{enumerate}

%     Respond with only the applicable category name(s), separated by commas. (e.g., ``WRONG\_LOGIC, THOUGHT\_SKIPPING'').
%     \end{boxK}

%     \caption{Prompt for classifying the failure mode. The failure mode is a multi-choice categorization. \junyuan{to update}}
%     \label{fig:failure_clf_prompt}
% \end{figure}